\documentclass[11pt,a4paper]{article}

\usepackage[margin=2.5cm]{geometry}
\usepackage{amsmath,amssymb,amsthm}
\usepackage{hyperref}
\usepackage{enumitem}
\usepackage{array}
\usepackage{booktabs}
\usepackage{pdflscape}

\theoremstyle{definition}
\newtheorem{definition}{Definition}[section]
\newtheorem{example}[definition]{Example}
\newtheorem{remark}[definition]{Remark}

\theoremstyle{plain}
\newtheorem{theorem}[definition]{Theorem}
\newtheorem{proposition}[definition]{Proposition}
\newtheorem{corollary}[definition]{Corollary}
\newtheorem{lemma}[definition]{Lemma}
\newtheorem{observation}[definition]{Observation}

\newcommand{\R}{\mathbb{R}}
\newcommand{\Z}{\mathbb{Z}}
\newcommand{\N}{\mathbb{N}}
\newcommand{\C}{\mathbb{C}}
\newcommand{\He}{\mathrm{He}}
\newcommand{\dfact}[1]{#1!!}

\title{\textbf{The Hermite Matching Matrix:\\
A Combinatorial Change of Basis\\
Between Hermite Polynomials and Hermite Functions}}
\author{Robin Langer and Claude Code}
\date{August 2026}

\begin{document}
\maketitle

\begin{abstract}
The Hermite polynomials $\{H_n(t)\}_{n \ge 0}$ are orthogonal with
respect to the Gaussian weight $e^{-t^2}$. The Hermite functions
$\psi_m(t) = c_m H_m(t)\, e^{-t^2/2}$ are obtained by absorbing the
square root of this weight into the polynomial; they form an orthonormal
basis for $L^2(\mathbb{R})$ and are eigenfunctions of the Fourier
transform. The $L^2$-projection of the $m$-th Hermite function onto the
Hermite polynomial basis---expressing $\psi_m$ as $\sum_k a_{mk} H_k$---defines
an infinite matrix $(a_{mk})$. We prove that, up to a single universal
constant, every entry of this matrix is an integer, and we conjecture an
explicit combinatorial formula: the rescaled entry $D_{mk}$ counts
\emph{signed bipartite matchings} of $m+k$ vertices, where cross-pairs
carry weight $+4$ and self-pairs carry weight $-2$. We give closed forms
for the edge entries $D_{0,2n} = (-1)^n\,(2n)!/n!$ and show that the
diagonal entries, divided by~$n!$, are integers yielding
the sequence
$\sum_k \binom{n}{2k}\binom{2k}{k}\,4^{n-2k}$ (OEIS A081671), with ordinary
generating function $1/\sqrt{1 - 8x + 12x^2}$ and exponential generating
function $e^{4x} I_0(2x)$, where $I_0$ is the modified Bessel function of
the first kind.
\end{abstract}

%% ===================================================================
\section{Setup}
%% ===================================================================

Let $H_n(t)$ denote the physicist's Hermite polynomials, defined by the
recurrence
\[
  H_{n+1}(t) = 2t\, H_n(t) - 2n\, H_{n-1}(t),
  \qquad H_0(t) = 1,\quad H_1(t) = 2t,
\]
and satisfying the orthogonality relation
\begin{equation}\label{eq:hermite-orth}
  \int_{-\infty}^{\infty} H_m(t)\, H_k(t)\, e^{-t^2}\, dt
  = 2^k\, k!\, \sqrt{\pi}\; \delta_{mk}.
\end{equation}
The \emph{Hermite functions} are
\begin{equation}\label{eq:hermite-fn}
  \psi_m(t) = c_m\, H_m(t)\, e^{-t^2/2},
  \qquad c_m = (2^m\, m!\, \sqrt{\pi})^{-1/2}.
\end{equation}
These form an orthonormal basis for $L^2(\R)$ and are eigenfunctions of
the Fourier transform with eigenvalues $(-i)^m$. Each $\psi_m$ is
entire, being a polynomial times a Gaussian.

The Hermite polynomials $\{H_k\}_{k \ge 0}$ form an orthogonal basis
for $L^2(\R, e^{-t^2}\,dt)$. The best $L^2$-approximation of $\psi_m$
by a polynomial of degree at most $N$ in this basis is
\[
  \psi_m^{(N)}(t) = \sum_{k=0}^{N} a_{mk}\, H_k(t),
\]
where
\begin{equation}\label{eq:amk}
  a_{mk}
  = \frac{c_m}{2^k\, k!\,\sqrt{\pi}}
    \int_{-\infty}^{\infty} H_m(t)\, H_k(t)\, e^{-3t^2/2}\, dt.
\end{equation}
The factor $e^{-3t^2/2}$ arises from the product of $e^{-t^2/2}$
(from~$\psi_m$) and $e^{-t^2}$ (from the Gaussian weight defining the
inner product). Our object of study is the infinite matrix $(a_{mk})$.

%% ===================================================================
\section{The Generating Function}
%% ===================================================================

\begin{proposition}\label{prop:genfn}
The integrals $I_{mk} = \int_{-\infty}^{\infty} H_m(t)\, H_k(t)\,
e^{-3t^2/2}\, dt$ have the generating function
\[
  \sum_{m,k \ge 0} \frac{s^m\, u^k}{m!\, k!}\; I_{mk}
  = \sqrt{\frac{2\pi}{3}}\;
    \exp\!\Bigl(-\frac{s^2 - 4su + u^2}{3}\Bigr).
\]
\end{proposition}

\begin{proof}
Using the generating function for Hermite polynomials,
$\sum_{n \ge 0} H_n(t)\, z^n / n! = e^{2tz - z^2}$, we compute
\begin{align*}
  \sum_{m,k} \frac{s^m u^k}{m!\,k!}\, I_{mk}
  &= \int_{-\infty}^{\infty}
     e^{2ts - s^2}\, e^{2tu - u^2}\, e^{-3t^2/2}\, dt \\
  &= \int_{-\infty}^{\infty}
     \exp\!\bigl(-\tfrac{3}{2}t^2 + 2(s+u)t - s^2 - u^2\bigr)\, dt \\
  &= \sqrt{\frac{2\pi}{3}}\;
     \exp\!\Bigl(\frac{2(s+u)^2}{3} - s^2 - u^2\Bigr) \\
  &= \sqrt{\frac{2\pi}{3}}\;
     \exp\!\Bigl(-\frac{s^2 - 4su + u^2}{3}\Bigr).
     \qedhere
\end{align*}
\end{proof}

\begin{corollary}\label{cor:parity}
$I_{mk} = 0$ whenever $m + k$ is odd (parity).
\end{corollary}
\begin{proof}
The exponent $-(s^2 - 4su + u^2)/3$ contains only terms of even total
degree in $(s,u)$.
\end{proof}

The matrix $(a_{mk})$ factors as
\begin{equation}\label{eq:factor}
  a_{mk} = R_{mk} \cdot \frac{\sqrt{2/3}}{\pi^{1/4}},
\end{equation}
where $R_{mk}$ is algebraic (involving only rationals and square roots),
and $\sqrt{2/3}/\pi^{1/4}$ is a universal constant independent of $m$
and~$k$.

%% ===================================================================
\section{The Rescaled Integer Matrix}
%% ===================================================================

Extracting coefficients from Proposition~\ref{prop:genfn} yields the
following.

\begin{proposition}\label{prop:Dmk}
Define
\begin{equation}\label{eq:Dmk}
  D_{mk} = 3^{(m+k)/2} \cdot \frac{I_{mk}}{\sqrt{2\pi/3}}.
\end{equation}
Then $D_{mk}$ is an integer for all $m, k \ge 0$ with $m + k$ even
(and zero when $m + k$ is odd).
\end{proposition}

\begin{proof}
Write
\[
  \exp\!\Bigl(-\frac{s^2 - 4su + u^2}{3}\Bigr)
  = \exp\!\Bigl(-\frac{s^2}{3}\Bigr)\;
    \exp\!\Bigl(\frac{4su}{3}\Bigr)\;
    \exp\!\Bigl(-\frac{u^2}{3}\Bigr).
\]
The coefficient of $s^m u^k$ is a sum over triples $(a, b, c)$ with
$m = 2a + b$ and $k = b + 2c\,$:
\[
  [s^m u^k] = \sum_{a,b,c}
  \frac{(-1/3)^a}{a!}\;\frac{(4/3)^b}{b!}\;\frac{(-1/3)^c}{c!}.
\]
Since $I_{mk} = m!\,k!\,[s^m u^k] \cdot \sqrt{2\pi/3}$, we have
$D_{mk} = 3^{(m+k)/2} \cdot m!\,k! \cdot [s^m u^k]$. Writing
$a + c = (m + k - 2b)/2$ and noting $3^{(m+k)/2}$ absorbs all powers
of~$3$ from the denominators while $4^b = 2^{2b}$ and $(-1)^{a+c}$
provide integer factors, every term is an integer.
\end{proof}

%% ===================================================================
\section{Apparent Matching Interpretation}
%% ===================================================================

The integers $D_{mk}$ appear to count signed bipartite matchings. We
have verified the following against the proved coefficient-extraction
formula of Proposition~\ref{prop:Dmk} for every pair $(m,k)$ with
$0 \le m, k \le 20$ (441~entries, zero discrepancies; see
Appendix~\ref{app:tables}).

\begin{observation}\label{obs:matching}
For all computed cases, $D_{mk}$ equals the sum over all perfect
matchings of $m + k$ vertices---$m$ of type~$A$ and $k$ of
type~$B$---where each matched pair contributes:
\begin{center}
\small
\begin{tabular}{lcc}
  \toprule
  Pair type & Raw weight & Rescaled weight \\
  \midrule
  Cross ($A$--$B$) & $4/3$ & $+4$ \\
  Self ($A$--$A$ or $B$--$B$) & $-2/3$ & $-2$ \\
  \bottomrule
\end{tabular}
\end{center}
Rescaled weights are obtained by multiplying through by $3^{(m+k)/2}$.
In closed form, this matching count would give
\begin{equation}\label{eq:Dmk-matching}
  D_{mk} = \sum_{\substack{c \ge 0 \\ m - c \text{ even} \\
  k - c \text{ even}}}
  \binom{m}{c}\binom{k}{c}\, c!\;
  (m - c - 1)!!\; (k - c - 1)!!\; 4^{c}\, (-2)^{(m+k-2c)/2},
\end{equation}
where $(2r-1)!! = 1 \cdot 3 \cdot 5 \cdots (2r-1)$ is the double
factorial and $(-1)!! = 1$ by convention. The factors count:
$\binom{m}{c}\binom{k}{c}$ ways to choose which vertices cross-pair,
$c!$~ways to match them (each cross-pair weighted $+4$), and
$(m - c - 1)!!\,(k - c - 1)!!$ self-matchings of the remainder
(each self-pair weighted $-2$).
\end{observation}

\begin{remark}
The matching interpretation is consistent with Wick's
theorem applied to the expectation
$\mathbb{E}[H_m(X)\,H_k(X)]$ for $X \sim N(0, 1/3)$.
A proof would follow from Isserlis' theorem at non-unit
variance.
\end{remark}

%% ===================================================================
\section{The Diagonal}
%% ===================================================================

\begin{theorem}\label{thm:diagonal}
The diagonal entries satisfy
\begin{equation}\label{eq:diag}
  \frac{D_{nn}}{n!} = \sum_{k=0}^{\lfloor n/2 \rfloor}
  \binom{n}{2k}\binom{2k}{k}\, 4^{n-2k}.
\end{equation}
This is OEIS sequence \textnormal{A081671}, with ordinary generating
function
\begin{equation}\label{eq:ogf}
  \sum_{n \ge 0} \frac{D_{nn}}{n!}\, x^n
  = \frac{1}{\sqrt{1 - 8x + 12x^2}}
  = \frac{1}{\sqrt{(1 - 2x)(1 - 6x)}},
\end{equation}
and exponential generating function
\begin{equation}\label{eq:egf}
  \sum_{n \ge 0} \frac{D_{nn}}{(n!)^2}\, x^n
  = e^{4x}\, I_0(2x),
\end{equation}
where $I_0$ is the modified Bessel function of the first kind.
\end{theorem}

\begin{proof}
Setting $m = k = n$ in~\eqref{eq:Dmk} and writing $j = (n-c)/2$ for
the number of self-pairs per side:
\[
  D_{nn} = \sum_{j=0}^{\lfloor n/2 \rfloor}
  \binom{n}{n-2j}^2 (n-2j)!\; \bigl((2j-1)!!\bigr)^2\;
  4^{n-2j}\, (-2)^{2j}.
\]
Since $(-2)^{2j} = 4^j$ and $(2j-1)!! = (2j)!/(2^j\, j!)$, the
summand becomes
\[
  \binom{n}{2j}^2 (n-2j)!\; \frac{(2j)!^2}{4^j\, (j!)^2}\; 4^j\, 4^{n-2j}
  = \binom{n}{2j}^2 (n-2j)!\; \frac{(2j)!^2}{(j!)^2}\; 4^{n-2j}.
\]
Dividing by $n!$ and simplifying:
\[
  \frac{D_{nn}}{n!}
  = \sum_{j} \binom{n}{2j}\; \frac{(2j)!}{(j!)^2}\; 4^{n-2j}
  = \sum_{j} \binom{n}{2j}\binom{2j}{j}\, 4^{n-2j}.
\]
For the OGF: let $a_n = \sum_k \binom{n}{2k}\binom{2k}{k}\,4^{n-2k}$.
Using $\sum_{n \ge 2k} \binom{n}{2k} (4x)^n
= (4x)^{2k}/(1-4x)^{2k+1}$,
\[
  \sum_{n} a_n x^n
  = \sum_k \binom{2k}{k}\, 4^{-2k}
    \frac{(4x)^{2k}}{(1-4x)^{2k+1}}
  = \frac{1}{1-4x} \sum_k \binom{2k}{k}
    \!\Bigl(\frac{x}{1-4x}\Bigr)^{\!2k}\!.
\]
Using the identity $\sum_k \binom{2k}{k} z^{2k} = 1/\sqrt{1-4z^2}$
with $z = x/(1-4x)$:
\[
  = \frac{1}{1-4x}\; \frac{1}{\sqrt{1 - 4x^2/(1-4x)^2}}
  = \frac{1}{\sqrt{(1-4x)^2 - 4x^2}}
  = \frac{1}{\sqrt{1 - 8x + 12x^2}}.
\]
For the EGF: $I_0(2x) = \sum_k x^{2k}/(k!)^2$, so
$e^{4x} I_0(2x) = \sum_n \bigl(\sum_k
\binom{n}{2k}\binom{2k}{k}\,4^{n-2k}/n!\bigr) x^n$
by Cauchy product, confirming~\eqref{eq:egf}.
\end{proof}

\begin{corollary}[Non-negativity]
The diagonal entries $D_{nn}/n!$ are non-negative integers for all
$n \ge 0$, being manifestly a sum of products of binomial coefficients.
\end{corollary}

The first several values of $D_{nn}/n!$ are:
\[
  1,\; 4,\; 18,\; 88,\; 454,\; 2424,\; 13236,\; 73392,\;
  411462,\; 2325976,\; \ldots
\]

%% ===================================================================
\section{The Edge Entries}
%% ===================================================================

The boundary entries of the matrix have closed forms.

\begin{proposition}\label{prop:edges}
For $k = 2n \ge 0$,
\[
  D_{0,2n} = (-1)^n\, \frac{(2n)!}{n!} = (-1)^n\, 2^n\, (2n-1)!!
\]
These are, up to sign, the \emph{rising factorials} $(2n)!/n!$\,.
\end{proposition}

\begin{proof}
When $m = 0$, the sum~\eqref{eq:Dmk} (assuming its validity) has only
the term $c = 0$, giving $D_{0,2n} = (2n-1)!!\,(-2)^n =
(-1)^n(2n)!/n!$. Alternatively, direct extraction from the generating
function: $[u^{2n}]\, e^{-u^2/3} = (-1/3)^n/n!$, so
$I_{0,2n} = \sqrt{2\pi/3}\,(2n)!\,(-1/3)^n/n!$ and
$D_{0,2n} = 3^n \cdot (2n)!\,(-1/3)^n/n! = (-1)^n(2n)!/n!$.
\end{proof}

%% ===================================================================
\section{Computed Values}
%% ===================================================================

The matrix $D_{mk}$ is symmetric ($D_{mk} = D_{km}$), vanishes when
$m + k$ is odd, and has alternating signs along each row. Since the
even-indexed and odd-indexed entries decouple, we present them as two
$8 \times 8$ blocks in Tables~\ref{tab:even} and~\ref{tab:odd}
(128~entries total). Table~\ref{tab:diagonal} lists $D_{nn}/n!$ for
$0 \le n \le 30$. All entries were computed directly from the proved
coefficient-extraction formula of Proposition~\ref{prop:Dmk}. The
matching formula~\eqref{eq:Dmk-matching} was independently checked
against these values for all $0 \le m, k \le 20$ (441~entries, zero
discrepancies).

%% ===================================================================
\section{Recovering the Original Matrix}
%% ===================================================================

The $L^2$-projection coefficients~\eqref{eq:amk} are recovered from the
integer matrix by
\[
  a_{mk} = \frac{D_{mk}}{3^{(m+k)/2}}
  \cdot \frac{1}{2^k\, k!\; \sqrt{2^m\, m!}}
  \cdot \frac{\sqrt{2/3}}{\pi^{1/4}}.
\]
The first factor undoes the rescaling by $3^{(m+k)/2}$, the second
accounts for the Hermite polynomial norms and the normalisation
constant~$c_m$, and the third is the universal constant.

%% ===================================================================
\section{Discussion}
%% ===================================================================

The change-of-basis matrix between Hermite functions and Hermite
polynomials is defined by a continuous, analytic process: $L^2$-projection
against a Gaussian weight. There is no \emph{a priori} reason for the
entries to be rational, let alone integer multiples of a single
constant.

The integrality of $D_{mk}$ (Proposition~\ref{prop:Dmk}) is proved by
direct coefficient extraction from the generating function. The matching
interpretation (Observation~\ref{obs:matching}) is verified
computationally for all $m + k \le 14$ but not proved here. A proof
would likely follow from Wick's theorem (Isserlis' theorem) applied to
the Gaussian expectation $\mathbb{E}[H_m(X)\,H_k(X)]$ for
$X \sim N(0, 1/3)$.

The non-negativity of the diagonal $D_{nn}/n!$ is independent of the
matching conjecture: it is a manifestly positive sum of products of
binomial coefficients (Theorem~\ref{thm:diagonal}), demanding a
bijective interpretation.

The appearance of the modified Bessel function $I_0$ in the diagonal's
EGF connects this matrix to wave equations and random walks. The OGF
$1/\sqrt{(1-2x)(1-6x)}$ is an algebraic function---the simplest class
beyond rational---with both branch points on the positive axis,
consistent with the positivity of the diagonal and its growth
rate~$6^n$. The matrix thus sits at an intersection of
combinatorics (matchings), special functions (Bessel), algebraic
generating functions, and classical analysis (Hermite, Fourier).

The matrix does not appear to have a standard name in the literature.
We propose \emph{Hermite matching matrix} for $(D_{mk})$ and note that
a complete proof of the matching interpretation, a characterisation of
the matrix's spectral properties, and its role in approximation theory
remain open.

%% ===================================================================
\appendix
\section{Tables}\label{app:tables}
%% ===================================================================

\begin{landscape}
\begin{table}[p]
\centering
\caption{The integer matrix $D_{mk}$ for even $m$ and $k$, $0 \le m,k \le 14$.}
\label{tab:even}
\scriptsize
\begin{tabular}{r|rrrrrrrr}
$m \!\backslash\! k$ & $0$ & $2$ & $4$ & $6$ & $8$ & $10$ & $12$ & $14$ \\
\hline
$0$ & $1$ & $-2$ & $12$ & $-120$ & $1680$ & $-30240$ & $665280$ & $-17297280$ \\
$2$ & $-2$ & $36$ & $-408$ & $6000$ & $-110880$ & $2479680$ & $-65197440$ & $1971889920$ \\
$4$ & $12$ & $-408$ & $10896$ & $-254880$ & $6471360$ & $-184222080$ & $5883736320$ & $-209435466240$ \\
$6$ & $-120$ & $6000$ & $-254880$ & $9529920$ & $-339534720$ & $12515731200$ & $-491854809600$ & $20808627840000$ \\
$8$ & $1680$ & $-110880$ & $6471360$ & $-339534720$ & $16590147840$ & $-792232358400$ & $38515400985600$ & $-1948173273446400$ \\
$10$ & $-30240$ & $2479680$ & $-184222080$ & $12515731200$ & $-792232358400$ & $48022189286400$ & $-2872542897100800$ & $173733606331084800$ \\
$12$ & $665280$ & $-65197440$ & $5883736320$ & $-491854809600$ & $38515400985600$ & $-2872542897100800$ & $208204121096294400$ & $-14960314003822387200$ \\
$14$ & $-17297280$ & $1971889920$ & $-209435466240$ & $20808627840000$ & $-1948173273446400$ & $173733606331084800$ & $-14960314003822387200$ & $1262574756680715878400$ \\
\end{tabular}
\end{table}

\begin{table}[p]
\centering
\caption{The integer matrix $D_{mk}$ for odd $m$ and $k$, $1 \le m,k \le 15$.}
\label{tab:odd}
\scriptsize
\setlength{\tabcolsep}{3pt}
\begin{tabular}{r|rrrrrrrr}
$m \!\backslash\! k$ & $1$ & $3$ & $5$ & $7$ & $9$ & $11$ & $13$ & $15$ \\
\hline
$1$ & $4$ & $-24$ & $240$ & $-3360$ & $60480$ & $-1330560$ & $34594560$ & $-1037836800$ \\
$3$ & $-24$ & $528$ & $-9120$ & $181440$ & $-4233600$ & $114428160$ & $-3528645120$ & $122464742400$ \\
$5$ & $240$ & $-9120$ & $290880$ & $-8588160$ & $266837760$ & $-9021196800$ & $334183449600$ & $-13545845913600$ \\
$7$ & $-3360$ & $181440$ & $-8588160$ & $369895680$ & $-15425303040$ & $658946534400$ & $-29586651494400$ & $1411067821363200$ \\
$9$ & $60480$ & $-4233600$ & $266837760$ & $-15425303040$ & $844050170880$ & $-45401368320000$ & $2476187275161600$ & $-139467481548595200$ \\
$11$ & $-1330560$ & $114428160$ & $-9021196800$ & $658946534400$ & $-45401368320000$ & $3021003695001600$ & $-198895976152473600$ & $13213366010836992000$ \\
$13$ & $34594560$ & $-3528645120$ & $334183449600$ & $-29586651494400$ & $2476187275161600$ & $-198895976152473600$ & $15600117724666675200$ & $-1214739624489431040000$ \\
$15$ & $-1037836800$ & $122464742400$ & $-13545845913600$ & $1411067821363200$ & $-139467481548595200$ & $13213366010836992000$ & $-1214739624489431040000$ & $109767194886547021824000$ \\
\end{tabular}
\end{table}
\end{landscape}

\begin{table}[p]
\centering
\caption{The diagonal sequence $D_{nn}/n!$ for $0 \le n \le 30$.
This is OEIS A081671, with OGF $1/\sqrt{1-8x+12x^2}$.}
\label{tab:diagonal}
\begin{tabular}{r|r||r|r}
$n$ & $D_{nn}/n!$ & $n$ & $D_{nn}/n!$ \\
\hline
$0$ & $1$ & $16$ & $487610895942$ \\
$1$ & $4$ & $17$ & $2838118247064$ \\
$2$ & $18$ & $18$ & $16547996212044$ \\
$3$ & $88$ & $19$ & $96635257790352$ \\
$4$ & $454$ & $20$ & $565107853947444$ \\
$5$ & $2424$ & $21$ & $3308820294176016$ \\
$6$ & $13236$ & $22$ & $19395905063796312$ \\
$7$ & $73392$ & $23$ & $113814537122646432$ \\
$8$ & $411462$ & $24$ & $668494299227072796$ \\
$9$ & $2325976$ & $25$ & $3929851838287363824$ \\
$10$ & $13233628$ & $26$ & $23120826355480783896$ \\
$11$ & $75682512$ & $27$ & $136129681993195136032$ \\
$12$ & $434662684$ & $28$ & $802049367833112712312$ \\
$13$ & $2505229744$ & $29$ & $4728540783318487334496$ \\
$14$ & $14482673832$ & $30$ & $27894081495241326235216$ \\
$15$ & $83940771168$ & & \\
\end{tabular}
\end{table}

\end{document}
